\documentclass[10pt,a4paper]{amsart}
\usepackage[margin=19mm]{geometry}
\usepackage{amsmath,amssymb,mathtools}
\usepackage[scr=rsfs,cal=euler]{mathalfa}
\usepackage{xfrac}
\usepackage[unicode,hidelinks,bookmarks=false]{hyperref}
\newcommand{\ie}{i.e.\ }
\newcommand{\cf}{cf.\ }
\newcommand{\resp}{respectively\ }
\newcommand{\Subsection}{\subsection}
\providecommand{\nobreakdash}{\nobreak\hbox{-}}
\setlength{\emergencystretch}{2em}
\multlinegap0pt
\usepackage{amssymb}
\newtheorem{proposition}{Proposition}
\newcommand{\Sph}{\mathbb S^1}
\newcommand{\cK}{\mathcal K}
\newcommand{\norm}[1]{\left\lVert #1\right\rVert}
\newcommand{\one}{\mathbf 1}
\title{Endpoint Criteria for One-Dimensional Bilinear Rough Singular Integrals}
\author{Binwei Dan, Qingying Xue\(^{*}\)}
\date{Source: arXiv:2607.17207v1 (CC BY 4.0)}
\begin{document}
\maketitle

\noindent\emph{Source and licence.} Endpoint Criteria for One-Dimensional Bilinear Rough Singular Integrals. Version arXiv:2607.17207v1 (CC BY 4.0), distributed by arXiv under the Creative Commons Attribution 4.0 International licence, http://creativecommons.org/licenses/by/4.0/, as recorded in the arXiv licence metadata for this exact version and shown on the abstract page https://arxiv.org/abs/2607.17207v1. Full licence notice: \texttt{licenses/arxiv-2607.17207-CC-BY-4.0.txt}.\par\smallskip

\noindent\textit{Editorial scope.} Counted unit: the article's proposition prop:Klog-not-Orlicz (source0.tex lines 3721--3727). The article's complete original proof of it is delivered with it (source0.tex lines 3729--3903, one \texttt{proof} environment of 175 lines). No mathematics is written, repaired or re-derived: the statement and the proof are the article's own lines, in the article's own notation, and no retained line is substituted or re-flowed.  No further source-local context is needed: every cross-reference of every retained line resolves inside the two delivered spans.  Nothing was omitted from any retained span, so the package carries no omission record.  No retained line contains a citation, so the wrapper carries no bibliography and no bibliography entry.  No retained line contains a figure environment, an image-inclusion command or drawing code, so no image is delivered and the package declares no asset dependency.  Editorial formatting only, in addition: an AMS \texttt{amsart} preamble that restates the article's own declarations and macros used by the retained lines, namely the environments \texttt{proposition} on separate counters, together with the standard packages those lines need.  The retained environments print in this document as Proposition 1 in order of appearance, whereas the article prints the same items with the numbering of its own full document.  The complete native source of the article as distributed in the arXiv e-print for this exact version is bundled unchanged as \texttt{sources/205546/source0.tex}.
	\begin{proposition}\label{prop:Klog-not-Orlicz}
		Let \(\alpha>0\) and \(\beta>0\). Then
		\[
		\cK_{1/2,\beta}(\Sph)\not\subset
		L(\log L)^\alpha(\Sph).
		\]
	\end{proposition}

	\begin{proof}
		Fix \(a\) with \(\frac12<a<1\).  We first construct a nonnegative function in
		\(\cK_a(\Sph)\) that does not belong to \(L(\log L)^\alpha(\Sph)\), and then
		subtract its mean. It will then suffice to use the embedding
		\(\cK_a(\Sph)\subset \cK_{1/2,\beta}(\Sph)\).
		
		Identify \(\Sph\) with \([0,2\pi)\) equipped with its normalized angular
		measure.  For each integer \(k\ge1\), set
		\[
		h_k:=4^k,\qquad w_k:=4^{-k}k^{-\alpha-1},
		\]
		and choose
		\[
		N_k:=\left\lceil
		w_k\,4^{\frac{k}{1-a}}k^{\frac{2}{1-a}}
		\right\rceil.
		\]
		Partition \(\Sph\) into \(N_k\) arcs of equal normalized measure \(1/N_k\). At
		the center of each arc, choose a closed subinterval \(I_{k,j}\) of normalized
		measure
		\[
		\sigma(I_{k,j})=\delta_k:=\frac{w_k}{N_k},
		\qquad 1\le j\le N_k.
		\]
		Let
		\[
		E_k:=\bigcup_{j=1}^{N_k}I_{k,j},
		\qquad
		\Omega_0(\theta):=\sum_{k=1}^\infty h_k\one_{E_k}(\theta).
		\]
		Each \(E_k\) is Borel, and hence the nonnegative function \(\Omega_0\) is
		Borel measurable. Moreover, \(\sigma(E_k)=w_k\) and, by Tonelli's theorem,
		\[
		\int_{\Sph}\Omega_0(\theta)\,d\sigma(\theta)
		=\sum_{k=1}^\infty h_k\sigma(E_k)
		=\sum_{k=1}^\infty k^{-\alpha-1}<\infty.
		\]
		Thus \(\Omega_0\) is finite almost everywhere and belongs to \(L^1(\Sph)\).
		We redefine it to be zero on the Borel null set where the displayed series
		is infinite, and set
		\[
		\mu_0:=\int_{\Sph}\Omega_0\,d\sigma,
		\qquad
		\Omega:=\Omega_0-\mu_0.
		\]
		Then \(\Omega\) is Borel measurable and has mean zero.
		
		We first show that \(\Omega_0\notin L(\log L)^\alpha(\Sph)\). Put
		\[
		A_k:=E_k\setminus\bigcup_{m>k}E_m.
		\]
		The sets \(A_k\) are pairwise disjoint, and
		\[
		\sigma(A_k)\ge \sigma(E_k)-\sum_{m>k}\sigma(E_m)
		=w_k-\sum_{m>k}w_m.
		\]
		Since \(w_m=4^{-m}m^{-\alpha-1}\),
		\[
		\sum_{m>k}w_m
		\le (k+1)^{-\alpha-1}\sum_{m=k+1}^\infty 4^{-m}
		=\frac13 (k+1)^{-\alpha-1}4^{-k}
		\le \frac13 w_k.
		\]
		Hence \(\sigma(A_k)\ge \frac23w_k\). Almost everywhere on \(A_k\) one has
		\(\Omega_0(\theta)\ge h_k\), and therefore
		\[
		\begin{aligned}
			\int_{\Sph}\Omega_0(\theta)
			\bigl(\log(e+\Omega_0(\theta))\bigr)^\alpha\,d\sigma(\theta)
			&\ge
			\sum_{k=1}^\infty h_k\bigl(\log(e+h_k)\bigr)^\alpha\sigma(A_k)\\
			&\gtrsim
			\sum_{k=1}^\infty 4^k k^\alpha\,4^{-k}k^{-\alpha-1}
			=\sum_{k=1}^\infty \frac1k
			=\infty.
		\end{aligned}
		\]
		Moreover, the divergence persists at every Luxemburg scale. For each
		\(\lambda>0\), choose \(k_\lambda\) so that, for every
		\(k\ge k_\lambda\),
		\(\log(e+h_k/\lambda)\gtrsim_\lambda k\), and hence
		\[
		\int_{\Sph}\Phi_\alpha\!\left(\frac{\Omega_0}{\lambda}\right)d\sigma
		\ge
		\sum_{k\ge k_\lambda}
		\frac{h_k}{\lambda}
		\bigl(\log(e+h_k/\lambda)\bigr)^\alpha\sigma(A_k)
		=\infty.
		\]
		Thus \(\Omega_0\notin L(\log L)^\alpha(\Sph)\). Since constants belong to
		\(L(\log L)^\alpha(\Sph)\) and this Orlicz space is a vector space, the
		identity \(\Omega_0=\Omega+\mu_0\) implies that
		\(\Omega\notin L(\log L)^\alpha(\Sph)\) as well.
		
		It remains to verify the fractional directional condition. For fixed
		\(\xi\in\Sph\), define
		\[
		I_k(\xi):=\int_{E_k}\frac{d\sigma(\theta)}{|\theta\cdot\xi|^a}.
		\]
		Let \(P_\xi=\{\xi^\perp,-\xi^\perp\}\) denote the two zeros of
		\(\theta\mapsto\theta\cdot\xi\) on \(\Sph\). Call \(I_{k,j}\) bad if its
		ambient partition arc meets \(P_\xi\). Since each point of \(P_\xi\) belongs to
		at most two adjacent partition arcs, there are at most four bad intervals.
		Their total contribution is bounded by
		\[
		\sum_{I_{k,j}\,{\rm bad}}
		\int_{I_{k,j}}|\theta\cdot\xi|^{-a}\,d\sigma(\theta)
		\lesssim_a \delta_k^{1-a},
		\]
		since, near either zero, \(|\theta\cdot\xi|\) is comparable to the angular
		distance to that zero.
		
		Group the remaining good intervals according to their grid distance \(m\) from
		\(P_\xi\). For each fixed \(m\), the number of such intervals is uniformly
		bounded, and on each of them
		\(|\theta\cdot\xi|\gtrsim m/N_k\). Hence
		\[
		\sum_{I_{k,j}\,{\rm good}}
		\int_{I_{k,j}}|\theta\cdot\xi|^{-a}\,d\sigma(\theta)
		\lesssim
		\delta_k N_k^a
		\sum_{m=1}^{N_k}m^{-a}
		\lesssim_a \delta_kN_k=w_k,
		\]
		since \(0<a<1\). Consequently
		\[
		\sup_{\xi\in\Sph} I_k(\xi)
		\lesssim_a \delta_k^{1-a}+w_k.
		\]
		Using \(N_k\ge w_k4^{k/(1-a)}k^{2/(1-a)}\), we have
		\[
		\delta_k=\frac{w_k}{N_k}
		\le 4^{-\frac{k}{1-a}}k^{-\frac{2}{1-a}},
		\qquad
		\delta_k^{1-a}\le 4^{-k}k^{-2}.
		\]
		Therefore
		\[
		\begin{aligned}
			\sup_{\xi\in\Sph}
			\int_{\Sph}\frac{\Omega_0(\theta)}{|\theta\cdot\xi|^a}
			\,d\sigma(\theta)
			&\le
			\sum_{k=1}^\infty h_k\sup_{\xi\in\Sph}I_k(\xi)\\
			&\lesssim_a
			\sum_{k=1}^\infty 4^k(4^{-k}k^{-2}+4^{-k}k^{-\alpha-1})
			<\infty.
		\end{aligned}
		\]
		Thus \(\Omega_0\in\cK_a(\Sph)\).  Moreover, since \(a<1\), rotational invariance
		and local integrability at the two zeros give
		\[
		C_a:=\sup_{\xi\in\Sph}
		\int_{\Sph}|\theta\cdot\xi|^{-a}\,d\sigma(\theta)<\infty.
		\]
		Consequently,
		\[
		\norm{\Omega}_{\cK_a}
		\le \norm{\Omega_0}_{\cK_a}+\mu_0 C_a<\infty,
		\]
		so the centered function \(\Omega\) also lies in \(\cK_a(\Sph)\).
		
		Finally, since \(a>\frac12\), for \(0<t\le1\) one has
		\[
		t^{-1/2}\left(\log\frac1t\right)^\beta\le C_{a,\beta}t^{-a}.
		\]
		Applying this with \(t=|\theta\cdot\xi|\) gives
		\[
		\norm{\Omega}_{\cK_{1/2,\beta}}
		\le C_{a,\beta}\norm{\Omega}_{\cK_a}<\infty.
		\]
		Therefore \(\Omega\) has mean zero, belongs to
		\(\cK_{1/2,\beta}(\Sph)\), and does not belong to
		\(L(\log L)^\alpha(\Sph)\).
	\end{proof}

\end{document}
